\section{Introduction}
\label{sec:intro}

\textquote{Land retains a focal point in Kenya's history. It was the basis upon which the struggle for independence was waged} \parencite{ndunguReportComissionInquiry2004}.
Since the British arrived in Kenya in 1895, the country and its population experienced numerous counts of land dispossession, injustices, and conflicts over land ownership and use.
Land that was previously owned and used by all Kenyans as a community resource was taken away and allocated to European settlers.
After independence in 1963, the Kenyan government aimed to redistribute the land to landless Kenyans through the implementation of settlement schemes, a common strategy used in many post-colonial countries.

The quote above is taken from a 2004 report on injustices in land allocation, suggesting that the issues did not end with independence.
One particular issue, the environmental impact of settlement schemes, was widely disregarded, the programs being implemented with the focus on social and economic development of the country.
This changed, as data that can be used to measure the impact of this program is now available and environmental degradation of land is a pressing issue.

High-resolution satellite images, available from 1980 onwards, allow for the monitoring of land use and land cover changes over time.
Their main advantage lies in the objective and consistent access, important in a country like Kenya,
with a history of lacking transparency and accountability in land allocation especially while the country was under dictatorship.
Combined with a dataset of settlement schemes including their boundaries and other properties, created by \cite{boonePromisedLandSettlement2021}, this allows for a detailed analysis of the environmental impact.

The main research objective of this thesis is to evaluate the causal environmental impact of settlement schemes in Kenya.
To achieve this, the study focuses on the intersection of land redistribution and environmental degradation.
As the social and economic implications of these programs are well documented, the implementation introduced an environmental shock to the environment.
Furthermore, their location in critical ecosystems, such as the Mau Forest, generates a high level of interest in understanding the environmental consequences for these areas and the surrounding regions.
The dataset of program boundaries allows the use of a Spatial Regression Discontinuity Design (RDD) to estimate the causal effect of settlement schemes on vegetation cover right at the scheme borders.

Beyond the main research objective, the methodology also allows for a more detailed analysis of subgroups of settlement schemes by factors such as size, population density, and the year of establishment.
Second, by utilizing forest cover data, the study can also analyze the impact of settlement schemes on forests in the surrounding areas, to further understand the consequences for this important ecosystem.

The remainder of this thesis is structured as follows: \cref{sec:background} provides a historical overview and definitions to understand the context of settlement schemes in Kenya.
\cref{sec:literature} reviews the existing literature on settlement schemes, land tenure security, deforestation, and spillover effects.
\cref{sec:data} describes the data used in this study. In \cref{sec:methods}, the methodology is explained, and the results are presented in \cref{sec:results}.
\cref{sec:discussion} places the results in the context of existing literature and outlines the limitations, and \cref{sec:conclusion} concludes the thesis.